% ----- Consignes exo 55 ----- %
\begin{td-exo}[Un algorithme de branchement pour le \(3\)-\textsc{SAT}]\,\\ % 55 
	Le but de l'exercice est de fournir un algorithme de résolution du problème NP-complet \(3\)-\textsc{SAT} qui fonctionne en temps exponentiel, mais le plus rapide possible.

	On précise d'abord quelques notations et définitions.
	Une instance \(F\) de \(3\)-\textsc{SAT} est une formule booléenne sous forme normale conjonctive telle que chaque clause contie,,e au plus trois littéraux.
	On notera \(n\) le nombre de variables de \(F\) et \(m\) son nombre de clauses.
	Soit \(c_i = (l_1 \lor l_2 \lor l_3)\) une clause de \(F\).
	On définit trois formules obtenues par \emph{assignement partiel} selon \(c_i\).
	La formule \(F_1 = F | l_1 = \textsc{true}\) est obtenue en enlevant à \(F\) toutes les clauses contenant \(l_1\) (car elles sont satisfaites) et en enlevant le littéral \(\lnot l_1\) de toutes les clauses contenant \(\lnot l_1\) (car elles ne peuvent pas être satisfaites par \(l_1\)).
	De même, \(F_2 = F | l_1 = \textsc{false}, l_2 = \textsc{true}\) est obtenue en enlevant à \(F\) toutes les clauses contenant \(l_2\) ou \(\lnot l_1\) et en enlevant le littéral \(\lnot l_2\) ou \(l_1\) de toutes les clauses contenant \(\lnot l_2\) ou \(l_1\).
	Enfin, \(F_3 = F | l_1 = \textsc{false}, l_2 = \textsc{false}, l_3 = \textsc{true}\) est obtenue en enlevant à \(F\) toutes les clauses contenant \(l_3\) ou \(\lnot l_1\) ou \(\lnot l_2\) et en enlevant le littéral \(\lnot l_3\) ou \(l_1\) ou \(l_2\) de toutes les clauses contenant \(\lnot l_3\) ou \(l_1\) ou \(l_2\).

	\begin{enumerate}
		\item Rappeler un algorithme de résolution de \(3\)-\textsc{SAT} qui fonctionne en temps \(O(m \cdot 2^n)\).
		
		\item Soit \(F\) une instance de \(3\)-\textsc{SAT} et \(c_i = (l_{i,1} \lor l_{i,2} \lor l_{i,3})\) une clause de \(F\).
		Montrer que si une formule obtenue par assignement partiel selon \(c_i\) est vide alors \(F\) est satisfiable.
		Que dire si une formule obtenue par assignement partiel contient une clause vide?

		\item Montrer que \(F\) est satisfiable si et seulement si l'une des trois formules \(F_1\), \(F_2\) ou \(F_3\) est satisfiable.
	\end{enumerate}
\end{td-exo}

% ----- Solutions exo 55 ----- %
\iftoggle{showsolutions}{ 
	\begin{td-sol}[]\ % 55
		\begin{enumerate}
			\item On essaie toutes les permutations de toutes les variables pour chacune des clauses.
			La complexité est alors de \(O(m \cdot 2^n)\) où \(m\) est le nombre de clauses et \(n\) le nombre de variables.

			\item Si une formule obtenue par assignement partiel est vide alors toutes les clauses de \(F\) sont satisfaites, et donc \(F\) est satisfiable.
			Si une formule obtenue par assignement partiel contient une clause vide alors il existe une clause de \(F\) qui n'est pas satisfaite, et donc \(F\) n'est pas satisfiable.

			\item Si \(F\) est satisfiable alors il existe une assignation des variables qui satisfait toutes les clauses de \(F\).
			En particulier, il existe une assignation des variables qui satisfait la clause \(c_i\).
			Par conséquent, il existe une assignation des variables qui satisfait l'une des trois formules \(F_1\), \(F_2\) ou \(F_3\).
			Inversement, si l'une des trois formules \(F_1\), \(F_2\) ou \(F_3\) est satisfiable alors il existe une assignation des variables qui satisfait toutes les clauses de \(F\), et donc \(F\) est satisfiable.
		\end{enumerate}
	\end{td-sol}
}{}
